\section{阅读口径}

本章解释报告中的产品、计量和控制术语，便于阅读，不替代基金合同、招募说明书、交易规则、账户条款或税务意见。除非四项实施闸门均已通过，以下定义均不构成任何产品的买卖、持有或现金用途指令。

\footnotesize
\begin{longtable}{p{2.9cm}p{7.2cm}}
\toprule
\textbf{术语} & \textbf{本报告中的含义与边界} \\
\midrule
\endfirsthead
\toprule
\textbf{术语} & \textbf{本报告中的含义与边界} \\
\midrule
\endhead
\bottomrule
\endfoot

\multicolumn{2}{l}{\textbf{产品与持有结构}} \\
\addlinespace[0.25em]
公募基础设施REITs（REITs） & REIT是\textit{Real Estate Investment Trust}的国际通用简称。本报告中的“公募REITs”特指公开募集基础设施证券投资基金，并不泛指境外REIT、房地产开发商股票或其他高分红基金。按R01所列产品特征，80\%以上基金资产投向基础设施资产支持证券并持有其全部份额，基金再通过特殊目的载体取得基础设施项目的完全所有权或经营权利，采取封闭式运作；投资者持有的是基金份额。该结构不构成现金分配或二级市场回报保证。 \\

基础设施资产支持证券／特殊目的载体 & 公募REIT持有基础设施项目权益的法律与交易载体。本报告使用该术语说明持有结构，不把它等同于投资者对项目现金流、项目价值或退出价格的保证。 \\

封闭式运作 & 本报告中的基础设施基金按封闭式方式运作；符合条件并经交易所审核同意上市后，基金份额可在交易所交易。二级市场价格可能与披露的净资产值不同。具体存续、交易与终止安排以基金合同和公告为准。 \\

ETF & 交易型开放式基金。本报告的\texttt{511010}和\texttt{510880}均为ETF研究锚，属于市场交易资产而非“受保障流动性”。ETF的历史分红、指数股息率或净值表现不自动等于未来实际到账现金，也不等于本策略的总回报。 \\

开放研究宇宙（T1） & 在授权、账户、税务、结算与流动性边界内，可被提名研究的直接股票、ETF、公募REIT、债券与现金工具集合。它不是持仓名单、推荐清单或自动准入名单；每项默认保持\(q=0\)。 \\

决策宇宙（T2） & 为一次具体比较而冻结的有限候选集。每一项必须有代码或身份、研究问题、数据截止日、责任人和证据路径。开放宇宙保证选择自由，T2保证比较可审计。 \\

批准可投池（T3）／批准组合（T4） & T3是完成结构适配的R0--R4研究并经独立证据、估值、风险和运营复核的具体标的；进入T3只代表可参与组合评审。T4是从T3按边际总回报、估值、风险、重叠、容量和流动性形成的获批组合。两者均不代表存在必须填满的资产类别配额。 \\

直接股票尽调 & 对单一发行人开展公司、行业、财务、盈利桥、治理、现金流/股利、当前估值、催化/证伪、A股流动性和组合重叠研究。它与ETF/REIT尽调同等严格，但证据模块因产品结构不同而不同。 \\

\addlinespace[0.25em]
\multicolumn{2}{l}{\textbf{现金与业绩计量}} \\
\addlinespace[0.25em]
可供分配金额 & 基础设施REIT判断现金分配基础的金额。报告仅采用R01所列规则口径：该金额在净利润基础上调整，并综合考虑项目公司持续发展、偿债能力和经营现金流。它不是净利润本身、不是持有人实际到账净现金，也不能直接当作按买入成本计算的收益率。 \\

现金分配／已结算现金记录 & 基金向持有人支付的现金。报告记录已结算、可核对的税前金额、税费、到账费用和净额，以支持总回报归因与产品尽调；卖出份额所得、未实现收益、净值上涨、外部入金和未结算应收款不得被误分类为分配。可能含返还性成分的分配须单列标记，并依预先声明的台账口径处理，不得称为收益或在总回报之外再加一次。 \\

TTM & \textit{Trailing Twelve Months}，即滚动过去12个月。它可用于查看产品历史分配、经营或估值数据，但不是下一年度分配或总回报预测，也不是本策略的独立验收目标。 \\

总回报 & 本报告将市价／NAV变动、现金分配、税费、交易成本和再平衡影响放入同一可对账口径。现金分配只在该口径内处理一次，不是额外叠加的一层收益或独立业绩目标。 \\

TWRR & \textit{Time-Weighted Rate of Return}，即时间加权收益率。每次外部入金、取现或划转均切分一个子期间，用于评估策略本身，不让投资者资金进出时点扭曲策略回报。 \\

MWRR／XIRR & 资金加权收益率／按日期的内部收益率。它以实际入金、取现和期末价值及其日期计算，用于反映投资者自身的资金体验；不能替代TWRR评价策略。 \\

返还性成分 & 可能反映资本返还而非经营或投资收益的现金组成部分。报告要求将其单列并在经济性质得到独立确认前排除出投资收益归因；最终处理不得造成总回报重复计数。 \\

NAV／IOPV／折溢价 & NAV是基金最近披露的净资产价值；IOPV是盘中参考估值；折溢价是市场价格相对适用参考价值的偏离。三者均不等于可成交价格。本报告不将REIT最新披露的NAV视为实时IOPV；缺少同一时段的价格、参考值和流动性资料时，资格闸门不能通过。 \\

\addlinespace[0.25em]
\multicolumn{2}{l}{\textbf{实施与风险控制}} \\
\addlinespace[0.25em]
资格闸门\(q\)／组合闸门\(q_{portfolio}\) & 报告用于表达实施资格的控制值。\(q=0\)表示该资产层因资料缺失、过期或未通过条件而不得配置；\(q=1\)只表示该层完成规定核验，仍须服从组合风险预算和独立实施决定。当前所有市场资产及组合均为\(q=0\)。 \\

受保障流动性 & 本报告的回退类别，指同时满足账户、产品、期限和相关保护条件的现金或直接国债类安排。它承接未获资格的金额，但不等于整套组合保本，也不保证达到5\%总回报目标。 \\

当前状态／历史参考向量 & “当前状态”是本报告日的模型动作：所有市场资产为\(q=0\)，T3/T4为空，不推定任何账户层事实。旧的10\%／15\%／40\%／25\%／10\%向量仅保留为被否决的首轮研究设计，用于说明固定类别配比的风险缺陷；它不是交易指令或未来配额。 \\

风险预算／共同压力 & 用具体候选的压力输入、相关性、集中度与流动性进行的治理测试，用于决定是否可进入实施审议；不是历史回撤、VaR、概率预测或最大损失承诺。旧向量的已知资产层共同压力小计为-25.00\%，高于-18\%风险行动线，因此不能作为固定配比启用。 \\

外部现金流调整后回撤\(DD_t\) & 相对外部现金流调整后高水位的账户市值跌幅。-12\%和-18\%是本报告对新增风险资金的治理行动线，不是自动卖出、止损保证或损失上限。 \\

风险簇／独立项目 & 本报告的内部集中度控制。基金代码不同不自动代表独立；需同时检查项目、管理人／发起人、行业、地区、主要租户或承购方、融资和相关性。至少五个独立项目只是最低必要条件，并不等于监管规定的“买满五只即可分散”。 \\

B0／B1 & 正式回放所需的两个预先冻结比较基准：受保障流动性机会成本和T2的点时跨资产组合。当前两者均未形成有效可比结果，不能用其推定未来5\%回报的概率。 \\
\end{longtable}
\normalsize

\sourcenote{公募基础设施REIT的法定结构和分配机制依据证监会《公开募集基础设施证券投资基金指引（试行）》（R01）。其余术语用于说明本报告既有的业绩、现金台账、风险和资格控制口径；不新增任何收益、现金或实施结论。}
